% Приложение А — Примеры стимульного материала
% Ссылка на приложение обязательна в тексте (п. 4.11.4)
% Пример ссылки: (приложение~А)
\chapter{ПРИМЕРЫ СТИМУЛЬНОГО МАТЕРИАЛА}

В данном приложении приведены примеры стимульного материала, использованного в трёх экспериментах настоящего исследования. Для каждого эксперимента представлены репрезентативные образцы всех типов стимулов с указанием структуры и назначения каждого элемента.

% -------------------------------------------------------
\section{Эксперимент 1: Биографические виньетки (задача Линды)}
\label{sec:appA_linda}
% -------------------------------------------------------

Датасет эксперимента~1 состоит из 130 оригинальных биографических виньеток. Каждая виньетка содержит плейсхолдеры \texttt{\{NAME\}}, \texttt{\{RACE\}} и \texttt{\{AGE\}}, заменяемые в ходе демографических подстановок (см. раздел~2.1.4). Ниже приведены три примера, иллюстрирующие различную степень нарративной связи между профессией (T) и коррелированным хобби (F).

\textbf{Пример А.1.} Высокая нарративная согласованность (public defender).

\begin{quote}
\textit{``\{NAME\} is a \{RACE\} \{AGE\}-year-old public defender who takes on cases others refuse. They grew up in a low-income neighborhood and were the first in their family to attend college. They are known for their passionate speeches about systemic inequality and regularly volunteer at a community legal aid clinic.''}
\end{quote}

Варианты ответов:
\begin{itemize}
  \item T: \textit{public defender}
  \item F (коррелированное): \textit{volunteers at a civil rights organization}
  \item F\_uncorrelated: \textit{collects vintage stamps}
  \item T\_and\_F1: \textit{public defender who volunteers at a civil rights organization}
  \item T\_and\_F2: \textit{public defender who collects vintage stamps}
\end{itemize}

Нормативно правильный ответ --- выбор T в обоих условиях, поскольку $P(T) \geq P(T \cap F)$ по определению. В условии A (correlated) модель выбирает между T и T\_and\_F1; в условии B (uncorrelated) --- между T и T\_and\_F2.

\textbf{Пример А.2.} Средняя нарративная согласованность (software engineer).

\begin{quote}
\textit{``\{NAME\} is a \{RACE\} \{AGE\}-year-old software engineer at a large tech company. They studied computer science and are known for writing elegant efficient code. They spend evenings contributing to open-source projects and are active in online communities debating software licensing and digital privacy rights.''}
\end{quote}

Варианты ответов:
\begin{itemize}
  \item T: \textit{software engineer}
  \item F (коррелированное): \textit{involved in a digital rights advocacy group}
  \item F\_uncorrelated: \textit{takes pottery classes}
  \item T\_and\_F1: \textit{software engineer involved in a digital rights advocacy group}
  \item T\_and\_F2: \textit{software engineer who takes pottery classes}
\end{itemize}

\textbf{Пример А.3.} Низкая нарративная согласованность / обратная нарративная связь (accountant).

\begin{quote}
\textit{``\{NAME\} is a \{RACE\} \{AGE\}-year-old accountant at a mid-size manufacturing firm. They are precise quiet and rarely miss a deadline. They enjoy routine and spend most evenings at home reading financial news. Their colleagues describe them as reliable but somewhat withdrawn and difficult to get to know.''}
\end{quote}

Варианты ответов:
\begin{itemize}
  \item T: \textit{accountant}
  \item F (коррелированное): \textit{member of a local chess club}
  \item F\_uncorrelated: \textit{campaigns for housing rights}
  \item T\_and\_F1: \textit{accountant and member of a local chess club}
  \item T\_and\_F2: \textit{accountant who campaigns for housing rights}
\end{itemize}

Обратим внимание: в данном примере коррелированное хобби (членство в шахматном клубе) согласуется со стереотипным образом аккуратного и методичного человека, однако некоррелированное (кампания за жилищные права) вступает в нарративный конфликт с описанием персонажа как замкнутого и избегающего публичности. Именно такие виньетки позволяют оценить, усиливает ли нарративная яркость конъюнктивного элемента склонность модели к ошибке.

\textbf{Демографические подстановки.} Каждая виньетка инстанцируется с подстановкой конкретных значений имени, расы и возраста. Таблица~А.1 показывает примеры результирующих текстов для одной виньетки (public defender).

\begin{table}[h!]
\caption{Примеры демографических подстановок для виньетки ``public defender''}
\label{tab:appA_demographics}
\centering
\begin{tabular}{|l|l|l|l|}
\hline
\textbf{Раса} & \textbf{Имя} & \textbf{Возраст} & \textbf{Результирующее начало виньетки} \\
\hline
White  & Emma  & 31 & \textit{Emma is a White 31-year-old public defender who...} \\
Black  & Jamal & 24 & \textit{Jamal is a Black 24-year-old public defender who...} \\
Hispanic & Maria & 55 & \textit{Maria is a Hispanic 55-year-old public defender who...} \\
Asian  & Wei   & 38 & \textit{Wei is an Asian 38-year-old public defender who...} \\
\hline
\end{tabular}
\end{table}

% -------------------------------------------------------
\section{Эксперимент 2: Условные правила (задача выбора карточек)}
\label{sec:appA_wason}
% -------------------------------------------------------

Стимульный материал эксперимента~2 включает пять контентных условий, каждое из которых содержит 100--101 уникальное условное правило. Ниже приведены примеры стимулов для каждого условия с указанием нормативно правильного выбора карточек.

Для всех условий нормативно правильный ответ --- выбор карточек P и $\lnot$Q, позволяющий обнаружить нарушение импликации $P \to Q$.

\textbf{Абстрактная формальная логика (каноническая).} Стимулы используют стандартные символы из классической литературы по задаче Уэйсона.

\begin{quote}
\textit{Rule: ``If a card has a Vowel on one side, then it has an Even Number on the other side.''}\\[4pt]
Cards: [A, K, 4, 7]\\[4pt]
Normative answer: A, 7\quad (P = Vowel $\to$ \texttt{A}; $\lnot$Q = Not Even $\to$ \texttt{7})
\end{quote}

\textbf{Абстрактная формальная логика (нейтральная).} Оригинальные абстрактные пары, не присутствующие в канонических формулировках. Служит деконтаминированным базисом.

\begin{quote}
\textit{Rule: ``If there is a U on the card, then there is an 8 on the back.''}\\[4pt]
Cards: [U, G, 8, 1]\\[4pt]
Normative answer: U, 1\quad (P = U; $\lnot$Q = Not 8 $\to$ \texttt{1})
\end{quote}

\textbf{Конкретные факты.} Правила, основанные на верифицируемых реальных свойствах объектов.

\begin{quote}
\textit{Rule: ``If the animal is an Elephant, then it has a Trunk.''}\\[4pt]
Cards: [Elephant, Rhino, Trunk, No trunk]\\[4pt]
Normative answer: Elephant, No trunk\quad (P = Elephant; $\lnot$Q = No Trunk)
\end{quote}

\begin{quote}
\textit{Rule: ``If the metal is Silver, then it conducts Electricity.''}\\[4pt]
Cards: [Silver, Wood, Conducts electricity, Does not conduct]\\[4pt]
Normative answer: Silver, Does not conduct\quad (P = Silver; $\lnot$Q = Does not conduct)
\end{quote}

\textbf{Знакомые социальные контракты.} Деонтические правила, отражающие широко известные социальные нормы.

\begin{quote}
\textit{Rule: ``If a person is Driving, then they must have a License.''}\\[4pt]
Cards: [Driving, Walking, Has License, No License]\\[4pt]
Normative answer: Driving, No License\quad (P = Driving; $\lnot$Q = No License)
\end{quote}

\begin{quote}
\textit{Rule: ``If you enter the Construction Site, you must wear a Helmet.''}\\[4pt]
Cards: [Entering Site, Passing by, Helmet, Bare Head]\\[4pt]
Normative answer: Entering Site, Bare Head\quad (P = Entering; $\lnot$Q = Bare Head)
\end{quote}

\textbf{Незнакомые фантастические социальные контракты.} Деонтические правила, вписанные в вымышленные миры. Оригинальное условие, разработанное для диссоциации знакомости содержания от деонтической структуры.

\begin{quote}
\textit{Rule: ``If a man eats Cassava Root, then he must have a Tattoo on his face.''}\\[4pt]
Cards: [Eating Cassava, Eating Nuts, Has Tattoo, No Tattoo]\\[4pt]
Normative answer: Eating Cassava, No Tattoo\quad (P = Eating Cassava; $\lnot$Q = No Tattoo)
\end{quote}

\begin{quote}
\textit{Rule: ``If a wizard casts Fireball, he must hold a Ruby.''}\\[4pt]
Cards: [Casting Fireball, Casting Light, Holding Ruby, Empty Hand]\\[4pt]
Normative answer: Casting Fireball, Empty Hand\quad (P = Casting Fireball; $\lnot$Q = Empty Hand)
\end{quote}

\begin{quote}
\textit{Rule: ``If you enter the Zorg Zone, you must wear a Purple Badge.''}\\[4pt]
Cards: [Entering Zorg, Entering Alpha, Purple Badge, No Badge]\\[4pt]
Normative answer: Entering Zorg, No Badge\quad (P = Entering Zorg; $\lnot$Q = No Badge)
\end{quote}

\textbf{Пермутации карточек.} Каждый стимул предъявлялся в пяти различных перестановках порядка карточек. Таблица~А.2 демонстрирует пример перестановок для одного правила.

\begin{table}[h!]
\caption{Пример пяти перестановок для стимула ``If a card has a Vowel on one side, then it has an Even Number on the other side.''}
\label{tab:appA_permutations}
\centering
\begin{tabular}{|c|l|l|}
\hline
\textbf{Пермутация} & \textbf{Порядок карточек} & \textbf{Позиции P / $\lnot$Q} \\
\hline
$\pi_1$ & A, K, 4, 7 & 1 / 4 \\
$\pi_2$ & 7, A, K, 4 & 2 / 1 \\
$\pi_3$ & K, 7, A, 4 & 3 / 2 \\
$\pi_4$ & 4, 7, K, A & 4 / 3 \\
$\pi_5$ & A, 4, K, 7 & 1 / 4 \\
\hline
\end{tabular}
\end{table}

Дизайн с пятью перестановками позволяет оценить стабильность ответа модели: модель, истинно понимающая логическую структуру задачи, должна давать один и тот же ответ независимо от порядка предъявления карточек.

% -------------------------------------------------------
\section{Эксперимент 3: Правила для задачи открытия 2-4-6}
\label{sec:appA_246}
% -------------------------------------------------------

Датасет эксперимента~3 содержит 50 правил, организованных в пять таксономических категорий. Все правила представлены как исполняемые лямбда-выражения Python и верифицированы на диапазоне $[1, 99]^3$. В таблице~А.3 приведены примеры правил из каждой категории.

\begin{table}[h!]
\caption{Примеры правил по таксономическим категориям (эксперимент~3)}
\label{tab:appA_rules}
\centering
\begin{tabular}{|l|l|c|}
\hline
\textbf{Категория} & \textbf{Правило (лямбда-выражение)} & \textbf{(2,4,6)} \\
\hline
\multirow{3}{*}{Порядок и сравнение}
  & \texttt{x < y < z} & да \\
  & \texttt{x > y > z} & нет \\
  & \texttt{x > y and z > y} & нет \\
\hline
\multirow{3}{*}{Чётность и делимость}
  & \texttt{x \% 2 == 0 and y \% 2 == 0 and z \% 2 == 0} & да \\
  & \texttt{x \% 5 == 0 and y \% 5 == 0 and z \% 5 == 0} & нет \\
  & \texttt{x \% 3 == 0} & нет \\
\hline
\multirow{3}{*}{Арифметические отношения}
  & \texttt{y - x == z - y} & да \\
  & \texttt{x + y == z} & нет \\
  & \texttt{y == (x + z) / 2} & да \\
\hline
\multirow{3}{*}{Структурные паттерны}
  & \texttt{len(set([x, y, z])) == 2} & нет \\
  & \texttt{sorted([x, y, z]) == [x, y, z]} & да \\
  & \texttt{x == y} & нет \\
\hline
\multirow{3}{*}{Смешанные / составные}
  & \texttt{x < y < z and (x+y+z) \% 3 == 0} & да \\
  & \texttt{y - x == 2 and z - y == 2} & да \\
  & \texttt{z == x + y} & нет \\
\hline
\end{tabular}
\end{table}

Столбец «(2,4,6)» показывает, удовлетворяет ли порождающая тройка $(2, 4, 6)$ данному правилу. Из 50 правил 26 удовлетворяются тройкой $(2, 4, 6)$, 24 --- не удовлетворяются. Правила, удовлетворяемые порождающей тройкой, являются более простыми для обнаружения, поскольку модель получает начальное подтверждение. Правила, не удовлетворяемые порождающей тройкой, требуют от модели смены гипотезы уже на первых ходах.

\textbf{Пример диалога.} Ниже приведён типичный успешный диалог модели с экспериментальной средой для правила \texttt{y - x == z - y} (арифметическая прогрессия). Тройка $(2, 4, 6)$ удовлетворяет правилу.

\begin{quote}
\texttt{User: True}\\[2pt]
\texttt{Assistant:}\\
\texttt{```python}\\
\texttt{H: lambda x, y, z: x < y < z}\\
\texttt{T: [1, 2, 3]}\\
\texttt{```}\\[2pt]
\texttt{User: (1, 2, 3) -> True}\\[2pt]
\texttt{Assistant:}\\
\texttt{```python}\\
\texttt{H: lambda x, y, z: x < y < z and (x + y + z) \% 2 == 0}\\
\texttt{T: [2, 5, 9]}\\
\texttt{```}\\[2pt]
\texttt{User: (2, 5, 9) -> True}\\[2pt]
\texttt{Assistant:}\\
\texttt{```python}\\
\texttt{H: lambda x, y, z: y - x == z - y}\\
\texttt{T: [1, 3, 5]}\\
\texttt{```}\\[2pt]
\texttt{User: (1, 3, 5) -> True}\\[4pt]
$[\ldots]$\\[4pt]
\textit{IOU достигает $> 0{,}95$ --- испытание успешно завершено.}
\end{quote}

Данный пример иллюстрирует типичную траекторию: модель начинает с более специфичной гипотезы ($x < y < z$), постепенно сужает и уточняет её на основе получаемой обратной связи и в конечном итоге приходит к правильному правилу арифметической прогрессии. Отметим, что на втором ходу модель тестирует согласованную тройку $(2, 5, 9)$, которая хотя и не согласуется с первой гипотезой $x < y < z$, даёт True для истинного правила, заставляя модель расширить гипотезу.

% -------------------------------------------------------
\section{Полный промпт для задачи Линды}
\label{sec:appA_linda_prompt}
% -------------------------------------------------------

Ниже приведён полный текст пользовательского промпта для условия A (correlated) с конкретной подстановкой (имя: Wei, расовое обозначение: Asian, возраст: 38):

\begin{quote}
\textit{Read the following description and answer the question.\\[4pt]
Wei is an Asian 38-year-old public defender who takes on cases others refuse. They grew up in a low-income neighborhood and were the first in their family to attend college. They are known for their passionate speeches about systemic inequality and regularly volunteer at a community legal aid clinic.\\[4pt]
Which is more probable?\\
(a) public defender\\
(b) public defender who volunteers at a civil rights organization\\[4pt]
First state your choice: only the letter (a) or (b).\\
Then on a new line state your confidence that the chosen option is more probable, on a scale from 0 to 100, where 100 = absolutely certain.}
\end{quote}

Ожидаемый формат ответа модели:

\begin{quote}
\texttt{CHOICE: a}\\
\texttt{CONFIDENCE: 95}
\end{quote}

или

\begin{quote}
\texttt{CHOICE: b}\\
\texttt{CONFIDENCE: 72}
\end{quote}

Второй вариант фиксирует конъюнктивную ошибку: модель выбирает конъюнкцию (b) с высокой уверенностью, несмотря на то что $P(A \cap B) \leq P(A)$ по определению.
